\section{Phase preparation without a known fundamental cell}
\label{sec:normalization}
The completion defect uses the free area of the \emph{whole} quotient.
A probability conditional on a successful return, or on launching near
one selected channel, cannot replace it. We now specify a preparation
which needs neither the period vectors nor the unknown obstacle list.
At finite scale it is approximate, and its error will remain in every
subsequent statistical bound.

\subsection{A periodic gate and a laboratory launch}
A record gate is a measurable function of a launched trajectory, returning
its two intrinsic endpoint coordinates on the selected branch, or a
failure symbol. It is \emph{periodic} if translating a launch and its
trajectory by any $\lambda\in\Lambda$ leaves the output unchanged.
Thus it counts all translates of its selected local branch by the same
local rule. It supplies no integer copy label. The local gate, its
paired-window identity and its coherent arclength reversal are retained
sensor capabilities. A gate addressed to one physical copy in the plane
is not periodic and is not covered by this construction.

The apparatus draws $q$ uniformly in the known laboratory square
$\Omega_L=[-L,L]^2$ and tests only whether $q$ is free. Solid launches
are rejected before a trajectory is launched. On a free launch it draws
a uniform direction, follows the billiard, and retains every gate outcome,
including failure. The apparatus need not report a geometric map of its
rejected positions. This is rejection in the \emph{launch location}, not
conditioning on success of a collision itinerary. Both rejected launch
attempts and all accepted free preparations are counted below.

Assume known bounds $A\ge A_0>0$, $V:=\covol\Lambda\le V_1$, and
that some fundamental parallelogram has diameter at most $b$. Its vectors,
position and free subset are not known. Such bounds follow from the
bounded lattice presentation in Section~\ref{sec:uniform}; for instance
$V_1=A_1+r_0\pi D_0^2$ is a harmless upper bound. Write $\nu_L$
for the projection modulo $\Lambda$ of an accepted free launch and its
direction, and $\mu$ for normalized quotient Liouville measure. Total
variation is $\|\nu-\mu\|_{\rm TV}=\sup_E|\nu(E)-\mu(E)|$.

\begin{proposition}[Uniform cell-free approximation]\label{prop:box}
For $L\ge2b$,
\begin{equation}\label{eq:box-tv}
 \|\nu_L-\mu\|_{\rm TV}\le
 \tau_L:=\min\left\{1,\frac{8V_1b}{A_0L}\right\}.
\end{equation}
The same bound holds after every periodic measurable record gate, for
every observation time and every cell of an endpoint histogram. The
free-launch probability is at least $p_0=A_0/(4V_1)$. Producing $N$
accepted free preparations takes at most $N/p_0$ launch attempts in
expectation, and at most
\begin{equation}\label{eq:launch-cost}
 \left\lceil\frac{8\{N+\log(1/\alpha)\}}{p_0}\right\rceil
\end{equation}
attempts with probability at least $1-\alpha$.
\end{proposition}
\begin{proof}
Tile the plane by a half-open fundamental parallelogram of diameter at
most $b$. Let $U$ be the union of tiles wholly contained in $\Omega_L$.
Every point of $[-L+b,L-b]^2$ lies in such a tile. Consequently
\[
 |\Omega_L\setminus U|\le8Lb,\qquad |U|\ge4(L-b)^2.
\]
If $U$ contains $n$ tiles, its free part has area $nA=(A/V)|U|$.
Conditional on a free launch in $U$, projection is \emph{exactly}
uniform on the free quotient; no knowledge of this tiling is used by
the apparatus. The projected measure from the remaining free part is
arbitrary, but its mixture weight is at most
\[
 \frac{8Lb}{(A_0/V_1)4(L-b)^2}\le\frac{8V_1b}{A_0L}.
\]
This proves \eqref{eq:box-tv}. Independent uniform direction and a
measurable gate cannot increase total variation. The free area in $U$
also gives launch acceptance probability at least
$(A_0/V_1)(1-b/L)^2\ge p_0$.

The number of attempts to get $N$ acceptances is a sum of independent
geometric waiting times, giving the expectation. For $m$ attempts its
upper tail is the probability that a binomial variable of mean at least
$mp_0$ is below $N$. If $m$ is the ceiling in \eqref{eq:launch-cost},
then $N\le mp_0/8$. The multiplicative Chernoff bound makes this
probability at most $\exp(-mp_0(7/8)^2/2)\le\alpha$.
\end{proof}

This argument is a direct periodic boundary-layer estimate, not an
assumption of mixing or equidistribution along one billiard orbit. It
remains valid in the presence of infinite horizons. Unknown hidden
bodies alter the free-launch distribution itself, so their area is not
removed by any later aperture normalization.

\subsection{Propagation through a finite histogram}
Fix nested active endpoint rectangles as in Proposition~\ref{prop:local-stable}.
Use a square mesh of side $h$, including a surrounding collar of cells.
The grid is centered so simultaneous arclength reversal permutes cells.
Let $f$ be a true cell-normalized density with a common
$C^{2,\beta}$ bound, $0<\beta\le1$. Cell probabilities satisfy
$p_B\le C h^2$. The large-square experiment instead has probabilities
$p_{B,L}$ with $|p_{B,L}-p_B|\le\tau_L$.

\begin{lemma}[Histogram estimate with preparation bias]
\label{lem:biased-histogram}
For $J$ fixed genuine records, $N$ independent accepted free preparations
at each of the two histogram times, and failure probability $\alpha$,
put $r_N=\log(CJN/\alpha)/N$. For $r_N<1$ and bandwidths with
$h^{-2}\le N$, a local quadratic reconstruction satisfies, simultaneously
for all records and both densities,
\begin{equation}\label{eq:biased-rate}
 e_N(h,L):=\|\widehat f-f\|_{C^2}
 \le C\{h^\beta+\sqrt{r_N}\,h^{-3}
                       +(r_N+\tau_L)h^{-4}\}
\end{equation}
with probability at least $1-\alpha$. In particular, choosing
\begin{equation}\label{eq:box-bandwidth}
 h=r_N^{1/(2\beta+6)},\qquad
 L\ge C r_N^{-(\beta+4)/(2\beta+6)}
\end{equation}
gives $e_N\le C r_N^{\beta/(2\beta+6)}$. The constants depend on
the physical and preparation bounds, not on a known cell or a known list
of obstacle shapes.
\end{lemma}
\begin{proof}
For one cell, Bernstein's inequality and
$p_{B,L}\le Ch^2+\tau_L$ bound its sampling error, after a union bound,
by $C\{h\sqrt{r_N}+\sqrt{\tau_Lr_N}+r_N\}$. Adding its deterministic
bias and using $\sqrt{\tau_Lr_N}\le(\tau_L+r_N)/2$ gives
\begin{equation}\label{eq:cell-error}
 |\widehat p_B-p_B|\le C\{h\sqrt{r_N}+r_N+\tau_L\}.
\end{equation}
There are at most $CJh^{-2}\le CJN$ cell events, so the stated logarithm
suffices. Different cell indicators from one launch need not be independent;
independence is used only across the $N$ launches for each indicator.

For completeness, averages of $1,z,z^2$ on three unit intervals centered
at $-1,0,1$ give the invertible matrix with rows $(1,i,i^2+1/12)$.
Its tensor square recovers a coordinatewise quadratic polynomial from
nine cell averages. On scale $h$, Taylor remainder of order $2+\beta$
gives derivative error $Ch^{2+\beta-j}$ at order $j\le2$.
A probability error $\varepsilon$ gives average error $\varepsilon h^{-2}$,
and differentiation costs a further $h^{-j}$. Patch the local polynomials
with a bounded-overlap smooth partition of unity whose derivatives of
order $j$ are $O(h^{-j})$. Leibniz's rule gives the same orders on the
smaller rectangle. At $j=2$, substitution of \eqref{eq:cell-error}
proves \eqref{eq:biased-rate}. No $C^2$ bound on the raw difference
$f_L-f$ is assumed; total variation was first converted to cell errors.

Under \eqref{eq:box-bandwidth}, $\tau_L\le C h^{\beta+4}$.
The bias and square-root terms balance at $r_N^{\beta/(2\beta+6)}$;
the linear Bernstein term is smaller. This proves the last assertion.
\end{proof}

A fixed pilot can select the two active absolute times. Uniform geometry
gives $cd^2\le P(T_0+d)\le Cd^2$ on a common right collar and zero
probability before onset. A fixed fine grid on the known onset bracket,
a fixed positive threshold and a collar after the upper bracket endpoint
place a reported crossing within a small fixed positive distance of $T_0$.
For sufficiently small $\tau_L+\sqrt{r_N}$ the same remains true with
sampling error. Adding two fixed positive offsets then places both
histograms on a common strictly active rectangle with separated times.
The pilot only locates an active interval; it is not an arbitrarily accurate
onset estimator. Fresh preparations for the two selected times justify
conditioning on the pilot. There are a prior-dependent bounded number
of pilot times per record, and their failures can be included in the
constant in Lemma~\ref{lem:biased-histogram}.
